\documentclass[11pt]{article}

\usepackage[margin=1in]{geometry}
\usepackage{amsmath}
\usepackage{booktabs}
\usepackage{enumitem}
\usepackage{float}
\usepackage{graphicx}
\usepackage{hyperref}
\usepackage{longtable}
\usepackage{xcolor}

\hypersetup{
  colorlinks=true,
  linkcolor=blue,
  urlcolor=blue
}

\title{Physics Embedded Transformers: Module Documentation}
\author{Project Documentation}
\date{\today}

\begin{document}
\maketitle

\tableofcontents
\listoffigures
\newpage

\section{Overview}

This document describes the project modules created so far for the
\texttt{physics\_embedded\_transformers} repository. The current implementation
builds a deterministic physical preprocessing stack for a microfluidic
asymmetric-loop device:

\begin{enumerate}[label=\arabic*.]
  \item Configuration loading.
  \item Centerline geometry construction.
  \item Device-level physical region labeling.
  \item Droplet fractional occupancy.
  \item Baseline branch hydraulics.
  \item Droplet branch-velocity validation.
\end{enumerate}

The modules are intentionally scoped to deterministic preprocessing and
baseline physics. They do not implement training, transformer features,
hydraulic parameter fitting, droplet--droplet interaction corrections, or
junction resistance models.

\section{Repository Layout}

\begin{longtable}{@{}ll@{}}
\toprule
Path & Purpose \\
\midrule
\texttt{configs/devices/asymmetric\_loop\_h100.yml} & Device geometry, calibration, region, and hydraulics settings. \\
\texttt{configs/experiments/video\_2.yml} & Experiment metadata, phase flow rates, and input paths. \\
\texttt{src/config/loader.py} & Experiment-to-device YAML loader. \\
\texttt{src/geometry/centerlines.py} & Centerline schema inspection, ordering, orientation, geometry quantities. \\
\texttt{src/geometry/types.py} & Dataclasses for typed centerline and device geometry containers. \\
\texttt{src/geometry/regions.py} & Device-level integer region label map construction and validation. \\
\texttt{src/occupancy/ellipse.py} & Local axis-aligned ellipse rasterization. \\
\texttt{src/occupancy/calculator.py} & Droplet fractional occupancy calculation and summary statistics. \\
\texttt{src/hydraulics/baseline.py} & Baseline branch-hydraulics equations. \\
\texttt{scripts/build\_geometry.py} & CLI for centerline geometry artifacts. \\
\texttt{scripts/build\_device\_regions.py} & CLI for device region labels and overlay. \\
\texttt{scripts/compute\_droplet\_occupancy.py} & CLI for droplet occupancy outputs. \\
\texttt{scripts/compute\_baseline\_hydraulics.py} & CLI for frame-level baseline hydraulics. \\
\texttt{scripts/validate\_droplet\_branch\_velocity.py} & Standalone normalized branch-velocity validation analysis. \\
\texttt{tests/} & Focused unit tests for geometry, regions, occupancy, hydraulics, and validation. \\
\bottomrule
\end{longtable}

\section{Configuration Layer}

\subsection{Files}

\begin{itemize}
  \item \texttt{src/config/loader.py}
  \item \texttt{src/config/\_\_init\_\_.py}
\end{itemize}

\subsection{Main Function}

\begin{verbatim}
load_experiment_config(
    experiment_path: str | Path,
    configs_root: str | Path = "configs",
) -> dict
\end{verbatim}

\subsection{Behavior}

The loader reads an experiment YAML file, extracts
\texttt{experiment.device\_id}, resolves the linked device file as:

\begin{verbatim}
configs/devices/<device_id>.yml
\end{verbatim}

It then loads the device YAML and verifies that:

\begin{verbatim}
device.id == experiment.device_id
\end{verbatim}

The returned structure is:

\begin{verbatim}
{
  "experiment": <full experiment YAML mapping>,
  "device": <full device YAML mapping>
}
\end{verbatim}

\subsection{Validation}

The loader raises clear exceptions for:

\begin{itemize}
  \item missing experiment files;
  \item empty YAML files;
  \item malformed YAML files;
  \item missing \texttt{experiment.device\_id};
  \item missing linked device files;
  \item mismatched experiment and device IDs.
\end{itemize}

\section{Device and Experiment Configuration}

\subsection{Device: \texttt{asymmetric\_loop\_h100}}

The device YAML stores fixed physical geometry and calibration:

\begin{itemize}
  \item calibration: \(\mu\mathrm{m/px} = 4.0\);
  \item channel width: \(100~\mu\mathrm{m}\);
  \item channel height: \(100~\mu\mathrm{m}\);
  \item left branch length: \(1791~\mu\mathrm{m}\);
  \item right branch length: \(1491~\mu\mathrm{m}\);
  \item junction definitions: \(100~\mu\mathrm{m} \times 100~\mu\mathrm{m}\);
  \item region artifact paths;
  \item isolated-droplet hydraulic resistance ratio.
\end{itemize}

The hydraulics section stores the dissertation provenance:

\begin{verbatim}
hydraulics:
  isolated_droplet_resistance:
    representation: equivalent_length
    ratio_to_short_branch: 0.15
\end{verbatim}

\subsection{Experiment: \texttt{video\_2}}

The experiment YAML stores:

\begin{itemize}
  \item device link: \texttt{device\_id: asymmetric\_loop\_h100};
  \item raw video location;
  \item processed data paths;
  \item continuous-phase flow rate;
  \item dispersed-phase flow rate;
  \item droplet nominal diameter and confinement regime.
\end{itemize}

For the current experiment:

\[
Q_{\mathrm{continuous}} = 1950~\mu\mathrm{L/hr},
\qquad
Q_{\mathrm{dispersed}} = 10~\mu\mathrm{L/hr}.
\]

The total mixture throughput used by baseline hydraulics is:

\[
Q_{\mathrm{mixture}} =
Q_{\mathrm{continuous}} + Q_{\mathrm{dispersed}}
= 1960~\mu\mathrm{L/hr}.
\]

\section{Centerline Geometry Module}

\subsection{Files}

\begin{itemize}
  \item \texttt{src/geometry/centerlines.py}
  \item \texttt{src/geometry/types.py}
  \item \texttt{scripts/build\_geometry.py}
  \item \texttt{tests/test\_centerline\_geometry.py}
\end{itemize}

\subsection{Input}

The geometry builder uses:

\begin{verbatim}
data/processed/2/centerlines.csv
\end{verbatim}

The discovered schema is:

\begin{center}
\begin{tabular}{@{}ll@{}}
\toprule
Column & Meaning \\
\midrule
\texttt{x} & centerline x coordinate in pixels \\
\texttt{y} & centerline y coordinate in pixels \\
\texttt{channel} & branch label \\
\bottomrule
\end{tabular}
\end{center}

The expected branch labels are:

\begin{verbatim}
inlet, outlet, left, right
\end{verbatim}

\subsection{Ordering and Orientation}

Each branch is represented as an ordered polyline. If explicit ordering columns
exist, they are preferred. Otherwise, the module attempts path reconstruction
from point adjacency. If adjacency reconstruction is unsuitable for the real
CSV, row order is preserved as the implicit sequence and this choice is reported.

Branches are oriented according to the physical flow direction:

\begin{itemize}
  \item inlet: upstream boundary to upper junction;
  \item left branch: upper junction to lower junction;
  \item right branch: upper junction to lower junction;
  \item outlet: lower junction to downstream boundary.
\end{itemize}

\subsection{Arc Length}

Segment length is Euclidean in pixel coordinates:

\[
\Delta s_k =
\sqrt{
(x_{k+1} - x_k)^2 +
(y_{k+1} - y_k)^2
}.
\]

Cumulative arc length is:

\[
s_0 = 0,
\qquad
s_k = \sum_{i=0}^{k-1}\Delta s_i.
\]

Physical arc length uses the device calibration:

\[
s_{\mu m} = s_{\mathrm{px}} \cdot 4.0.
\]

\subsection{Tangents and Normals}

Interior tangents use centered finite differences. Endpoints use one-sided
finite differences. Tangents are normalized to unit length.

The normal convention is:

\[
n_x = -t_y,
\qquad
n_y = t_x.
\]

This normal is a consistent mathematical convention only; it is not labeled as
inward or outward.

\subsection{Typed Containers}

\texttt{BranchCenterline} stores:

\begin{itemize}
  \item branch name;
  \item ordered \(x,y\) coordinates;
  \item cumulative arc length in pixels and micrometres;
  \item tangent vectors;
  \item normal vectors;
  \item total length in pixels and micrometres.
\end{itemize}

\texttt{DeviceGeometry} stores:

\begin{itemize}
  \item device ID;
  \item calibration;
  \item branch centerline dictionary;
  \item upper and lower junction coordinates;
  \item endpoint mismatch diagnostics.
\end{itemize}

\subsection{Outputs}

\begin{verbatim}
outputs/geometry/asymmetric_loop_h100/centerline_geometry.csv
outputs/geometry/asymmetric_loop_h100/geometry_summary.json
outputs/geometry/asymmetric_loop_h100/device_geometry.npz
outputs/geometry/asymmetric_loop_h100/centerline_geometry.png
\end{verbatim}

\subsection{Execution}

\begin{verbatim}
.\.venv\Scripts\python.exe -m scripts.build_geometry \
  --experiment configs\experiments\video_2.yml \
  --overwrite
\end{verbatim}

\section{Device Region Labeling Module}

\subsection{Files}

\begin{itemize}
  \item \texttt{src/geometry/regions.py}
  \item \texttt{scripts/build\_device\_regions.py}
  \item \texttt{tests/test\_device\_regions.py}
\end{itemize}

\subsection{Purpose}

The region-labeling module creates a device-level integer label map for physical
regions. It is independent of a particular droplet frame and is intended as the
spatial support for downstream fractional occupancy calculations.

\subsection{Label IDs}

\begin{center}
\begin{tabular}{@{}cl@{}}
\toprule
ID & Region \\
\midrule
0 & unassigned/background \\
1 & inlet \\
2 & outlet \\
3 & left branch \\
4 & right branch \\
5 & upper junction \\
6 & lower junction \\
\bottomrule
\end{tabular}
\end{center}

\subsection{Construction}

The label map is constructed from:

\begin{itemize}
  \item the binary channel mask;
  \item original labeled centerlines;
  \item junction definitions from the device YAML.
\end{itemize}

Junction boxes are converted from micrometres to pixels using:

\[
\mathrm{size}_{px} =
\frac{\mathrm{size}_{\mu m}}{\mu m/px}.
\]

For the current device:

\[
100~\mu\mathrm{m} / 4~\mu\mathrm{m/px} = 25~\mathrm{px}.
\]

Pixels outside the junction boxes are assigned to the nearest compatible
centerline while remaining clipped to the channel mask.

\subsection{Validation}

The module validates:

\begin{itemize}
  \item label-map shape equals channel-mask shape;
  \item valid label IDs only;
  \item no nonzero labels outside the channel mask;
  \item nonempty physical regions;
  \item branch and inlet/outlet connectivity to the appropriate junctions.
\end{itemize}

\subsection{Outputs}

\begin{verbatim}
data/geometry/asymmetric_loop_h100/region_labels.npy
data/geometry/asymmetric_loop_h100/region_metadata.json
data/geometry/asymmetric_loop_h100/region_overlay.png
\end{verbatim}

\begin{figure}[H]
  \centering
  \includegraphics[width=0.95\linewidth,height=0.75\textheight,keepaspectratio]{../data/geometry/asymmetric_loop_h100/region_overlay.png}
  \caption{Device-level region overlay produced by \texttt{scripts/build\_device\_regions.py}. The overlay is generated from the precomputed integer region labels, not from visualization-time region inference.}
  \label{fig:region-overlay}
\end{figure}

\subsection{Execution}

\begin{verbatim}
.\.venv\Scripts\python.exe -m scripts.build_device_regions \
  --experiment configs\experiments\video_2.yml \
  --frame-index 0 \
  --overwrite
\end{verbatim}

\section{Droplet Fractional Occupancy Module}

\subsection{Files}

\begin{itemize}
  \item \texttt{src/occupancy/ellipse.py}
  \item \texttt{src/occupancy/calculator.py}
  \item \texttt{src/occupancy/\_\_init\_\_.py}
  \item \texttt{scripts/compute\_droplet\_occupancy.py}
  \item \texttt{tests/test\_droplet\_occupancy.py}
\end{itemize}

\subsection{Inputs}

The occupancy pipeline reads:

\begin{itemize}
  \item \texttt{outputs/physics/video\_2/droplet\_occupancy.csv} for downstream use after generation;
  \item \texttt{data/processed/2/tracked\_features.csv} as the raw tracking input;
  \item \texttt{data/geometry/asymmetric\_loop\_h100/region\_labels.npy};
  \item \texttt{data/geometry/asymmetric\_loop\_h100/region\_metadata.json};
  \item device calibration from the YAML.
\end{itemize}

The tracked-feature columns used are:

\begin{verbatim}
frame, track_id, centroid_x, centroid_y, bbox_w, bbox_h
\end{verbatim}

\subsection{Ellipse Footprint}

Each droplet footprint is approximated as an axis-aligned ellipse fitted to its
tracked bounding box:

\[
\left(\frac{x - x_c}{w/2}\right)^2
+
\left(\frac{y - y_c}{h/2}\right)^2
\leq 1.
\]

Rasterization uses pixel-center coordinates in the same image convention as the
region label map:

\begin{itemize}
  \item \(x\) increases with image columns;
  \item \(y\) increases with image rows.
\end{itemize}

Only the local clipped bounding window around each ellipse is rasterized.

\subsection{Raw Occupancy}

For physical region \(k\), raw occupancy is:

\[
W_{k,\mathrm{raw}} =
\frac{\#\{\mathrm{ellipse~pixels~with~label~}k\}}
{\#\{\mathrm{all~rasterized~ellipse~pixels~inside~the~image}\}}.
\]

The denominator is the complete rasterized ellipse area after image-boundary
clipping, not only the channel-overlapping part.

The raw physical coverage is:

\[
C_{\mathrm{phys}} =
\sum_{k=1}^{6} W_{k,\mathrm{raw}}.
\]

\subsection{Normalized Occupancy}

For every sample with:

\[
C_{\mathrm{phys}} > \epsilon,
\]

the normalized six-region occupancy vector is:

\[
W_k =
\frac{W_{k,\mathrm{raw}}}{C_{\mathrm{phys}}}.
\]

The normalized occupancy vector is the canonical representation for downstream
ML and physics modules:

\[
W =
\begin{bmatrix}
W_{\mathrm{inlet}} &
W_{\mathrm{outlet}} &
W_{\mathrm{left}} &
W_{\mathrm{right}} &
W_{\mathrm{upper~junction}} &
W_{\mathrm{lower~junction}}
\end{bmatrix}.
\]

For computable samples:

\[
\sum_k W_k = 1.
\]

\subsection{Diagnostic Low Coverage}

The threshold \(0.95\) is retained as a diagnostic flag:

\[
\mathrm{low\_physical\_coverage} =
C_{\mathrm{phys}} < 0.95.
\]

Low-coverage samples are still normalized when \(C_{\mathrm{phys}} > \epsilon\).
They are not called invalid merely because they fall below the diagnostic
threshold.

\subsection{CSV Columns}

Important output columns include:

\begin{itemize}
  \item diagnostic raw quantities:
  \texttt{w\_inlet\_raw},
  \texttt{w\_outlet\_raw},
  \texttt{w\_left\_raw},
  \texttt{w\_right\_raw},
  \texttt{w\_upper\_junction\_raw},
  \texttt{w\_lower\_junction\_raw},
  \texttt{physical\_coverage\_raw},
  \texttt{low\_physical\_coverage};
  \item pipeline quantities:
  \texttt{occupancy\_computable},
  \texttt{w\_inlet},
  \texttt{w\_outlet},
  \texttt{w\_left},
  \texttt{w\_right},
  \texttt{w\_upper\_junction},
  \texttt{w\_lower\_junction},
  \texttt{normalized\_occupancy\_sum};
  \item diagnostics:
  \texttt{image\_boundary\_clipped},
  \texttt{dominant\_region},
  \texttt{number\_active\_regions}.
\end{itemize}

No \texttt{w\_unassigned} component is created.

\subsection{Outputs}

\begin{verbatim}
outputs/physics/video_2/droplet_occupancy.csv
outputs/physics/video_2/droplet_occupancy_summary.json
outputs/physics/video_2/occupancy_diagnostics.png
outputs/physics/video_2/occupancy_spot_checks.png
\end{verbatim}

\subsection{Execution}

\begin{verbatim}
.\.venv\Scripts\python.exe -m scripts.compute_droplet_occupancy \
  --experiment configs\experiments\video_2.yml \
  --overwrite
\end{verbatim}

\section{Baseline Branch Hydraulics Module}

\subsection{Files}

\begin{itemize}
  \item \texttt{src/hydraulics/baseline.py}
  \item \texttt{src/hydraulics/\_\_init\_\_.py}
  \item \texttt{scripts/compute\_baseline\_hydraulics.py}
  \item \texttt{tests/test\_baseline\_hydraulics.py}
\end{itemize}

\subsection{Purpose}

The baseline hydraulics module computes a deterministic frame-level hydraulic
state using normalized branch occupancies and a simple equivalent-resistance
model. It deliberately excludes inlet, outlet, and junction occupancies from the
baseline resistance calculation.

\subsection{Effective Branch Occupancy}

For frame \(t\):

\[
N_{\mathrm{left,eff}}(t)
=
\sum_{d \in t} W_{\mathrm{left}}(d),
\]

\[
N_{\mathrm{right,eff}}(t)
=
\sum_{d \in t} W_{\mathrm{right}}(d).
\]

These are continuous effective occupancies, not integer droplet counts.

\subsection{Isolated-Droplet Equivalent Length}

The dissertation-derived isolated-droplet resistance length is:

\[
\ell_d = r_d L_{\mathrm{short}},
\]

where:

\[
r_d = 0.15.
\]

For the current device:

\[
L_{\mathrm{short}} = L_{\mathrm{right}} = 1491~\mu\mathrm{m},
\]

so:

\[
\ell_d = 0.15 \times 1491
= 223.65~\mu\mathrm{m}.
\]

\subsection{Effective Branch Lengths}

\[
L_{\mathrm{left,eff}}(t)
=
L_{\mathrm{left}}
+ \ell_d N_{\mathrm{left,eff}}(t),
\]

\[
L_{\mathrm{right,eff}}(t)
=
L_{\mathrm{right}}
+ \ell_d N_{\mathrm{right,eff}}(t).
\]

\subsection{Total Mixture Flow}

The hydraulic throughput is the total mixture flow:

\[
Q_{\mathrm{mix}}
=
Q_{\mathrm{continuous}}
+
Q_{\mathrm{dispersed}}.
\]

For \texttt{video\_2}:

\[
Q_{\mathrm{mix}}
= 1950 + 10
= 1960~\mu\mathrm{L/hr}.
\]

\subsection{Parallel Branch Split}

Because resistance is proportional to equivalent length and both branches have
the same cross-section, the flow split is:

\[
Q_{\mathrm{left}}
=
Q_{\mathrm{mix}}
\frac{L_{\mathrm{right,eff}}}
{L_{\mathrm{left,eff}} + L_{\mathrm{right,eff}}},
\]

\[
Q_{\mathrm{right}}
=
Q_{\mathrm{mix}}
\frac{L_{\mathrm{left,eff}}}
{L_{\mathrm{left,eff}} + L_{\mathrm{right,eff}}}.
\]

The reconstructed total is:

\[
Q_{\mathrm{reconstructed}}
=
Q_{\mathrm{left}} + Q_{\mathrm{right}}.
\]

The conservation diagnostic is:

\[
\Delta Q =
Q_{\mathrm{reconstructed}} - Q_{\mathrm{mix}}.
\]

\subsection{Superficial Mixture Velocity}

The branch velocity columns are total superficial mixture velocities:

\[
U =
\frac{Q}{W H}.
\]

The unit conversion is:

\[
1~\mu\mathrm{L}
=
10^9~\mu\mathrm{m}^3,
\qquad
1~\mathrm{hr}
=
3600~\mathrm{s}.
\]

Thus:

\[
U_{\mu m/s}
=
\frac{Q_{\mu L/hr} \cdot 10^9 / 3600}
{W_{\mu m}H_{\mu m}}.
\]

\subsection{Canonical ML Features}

The canonical hydraulic features intended for the future transformer are only:

\begin{verbatim}
left_velocity_um_s
right_velocity_um_s
\end{verbatim}

These names are kept stable for downstream compatibility, but their documented
meaning is:

\begin{itemize}
  \item \texttt{left\_velocity\_um\_s}: total superficial mixture velocity in the left branch;
  \item \texttt{right\_velocity\_um\_s}: total superficial mixture velocity in the right branch.
\end{itemize}

Flow-rate columns remain in the CSV for diagnostics and conservation checks, but
they are not intended as independent ML inputs because they are deterministically
related to velocity by the fixed channel cross-section.

\subsection{Baseline Assumptions}

\begin{enumerate}[label=\arabic*.]
  \item The total volumetric throughput is the sum of continuous- and dispersed-phase input flow rates.
  \item The total mixture flow is split between the two branches according to effective hydraulic resistances.
  \item Reported branch velocities are superficial mixture velocities based on the full channel cross-sectional area.
  \item No phase-specific slip or separate phase velocity is modeled.
  \item Branch resistance is proportional to equivalent branch length.
  \item Both branches have identical cross-sections.
  \item Fluid properties are identical between branches.
  \item Droplet resistance contributions are additive.
  \item Each isolated droplet contributes \(0.15\) times the short-branch resistance length.
  \item Fractional branch occupancy scales the droplet contribution linearly.
  \item Inlet, outlet, upper-junction, and lower-junction occupancies do not contribute to baseline resistance.
  \item Droplet--droplet interaction corrections are not modeled.
  \item Junction dynamics are reserved for the future learned model.
\end{enumerate}

\subsection{Outputs}

\begin{verbatim}
outputs/physics/video_2/baseline_hydraulic_state.csv
outputs/physics/video_2/baseline_hydraulic_summary.json
outputs/physics/video_2/baseline_hydraulic_diagnostics.png
\end{verbatim}

\subsection{Execution}

\begin{verbatim}
.\.venv\Scripts\python.exe -m scripts.compute_baseline_hydraulics \
  --experiment configs\experiments\video_2.yml \
  --overwrite
\end{verbatim}

\section{Droplet Branch-Velocity Validation Module}

\subsection{Files}

\begin{itemize}
  \item \texttt{scripts/validate\_droplet\_branch\_velocity.py}
  \item \texttt{tests/test\_droplet\_branch\_velocity\_validation.py}
\end{itemize}

\subsection{Purpose}

This standalone validation analysis compares measured droplet-speed variation
against the calculated baseline branch superficial-velocity variation. It does
not modify the tracking, occupancy, or hydraulic preprocessing layers.

The analysis is intentionally restricted to branch-interior samples. It excludes
inlet, outlet, junction, and branch-transition samples by requiring a normalized
branch occupancy of at least \(0.95\) in exactly one branch:

\[
w_{\mathrm{left}} \ge 0.95 \quad \mathrm{xor} \quad
w_{\mathrm{right}} \ge 0.95.
\]

\subsection{Measured Velocity Derivation}

For the current \texttt{video\_2} tracking table, the canonical file is:

\begin{verbatim}
data/processed/2/tracked_features.csv
\end{verbatim}

The table does not contain existing \texttt{vx} and \texttt{vy} columns, so the
validator derives velocity from the configured centroid columns:

\begin{verbatim}
centroid_x, centroid_y
\end{verbatim}

For each track, rows are sorted by frame. Interior samples receive centered
finite-difference velocity only when both neighboring observations are exactly
one frame apart:

\[
v_x(t) = \frac{x(t+1) - x(t-1)}{2},
\qquad
v_y(t) = \frac{y(t+1) - y(t-1)}{2}.
\]

The measured speed is:

\[
u_{\mathrm{drop}}(t) =
\sqrt{v_x(t)^2 + v_y(t)^2}.
\]

Endpoint samples and samples adjacent to tracking gaps are excluded. Track
boundaries are never mixed. If future tracking files provide finite
\texttt{vx}/\texttt{vy} columns, those columns are preferred and the metadata
records \texttt{tracked\_velocity\_source = existing\_velocity\_columns}.

\subsection{Unit Interpretation}

For the current analysis:

\begin{itemize}
  \item measured droplet speed units are \texttt{px/frame};
  \item calculated branch velocity units are \(\mu\mathrm{m}/\mathrm{s}\);
  \item \texttt{tracked\_velocity\_source} is
  \texttt{centered\_position\_difference}.
\end{itemize}

The validation does not compare absolute velocity magnitudes. Each signal is
standardized independently within each selected \((\texttt{track\_id},
\texttt{branch})\) segment. Therefore, the plots and correlations compare
normalized temporal co-variation only. No frame-rate or pixel-scale conversion
is required for this normalized comparison.

\subsection{Normalization}

For every selected track segment, measured speed and calculated branch velocity
are normalized separately using sample standard deviation with \(ddof=1\):

\[
z_{\mathrm{drop}}(t) =
\frac{u_{\mathrm{drop}}(t) - \bar{u}_{\mathrm{drop}}}
{\mathrm{std}(u_{\mathrm{drop}})},
\]

\[
z_{\mathrm{branch}}(t) =
\frac{u_{\mathrm{branch}}(t) - \bar{u}_{\mathrm{branch}}}
{\mathrm{std}(u_{\mathrm{branch}})}.
\]

Tracks with negligible measured-speed or branch-velocity variance are rejected.
No smoothing, lag optimization, or physical speed-scale fitting is performed.

\subsection{Selection}

The intended default minimum branch-interior sample count was 100. The current
\texttt{video\_2} branch-interior segments are shorter than this: left-branch
segments top out near 85 samples, while right-branch segments top out near 67
samples. To avoid silently excluding the entire right branch, this analysis uses:

\begin{verbatim}
minimum_branch_interior_samples: 60
\end{verbatim}

The cutoff is recorded in the summary JSON. Candidate ranking is deterministic:
within each branch, candidates are ranked by
\texttt{std\_branch\_velocity\_um\_s} descending, using \texttt{track\_id} as a
tie-breaker. Correlation is not used for selection.

For \texttt{video\_2}, eligible candidates were imbalanced:

\begin{center}
\begin{tabular}{@{}lr@{}}
\toprule
Quantity & Count \\
\midrule
Total eligible candidates & 3595 \\
Left-branch candidates & 1148 \\
Right-branch candidates & 2447 \\
Selected candidates & 10 \\
Pooled normalized samples & 742 \\
\bottomrule
\end{tabular}
\end{center}

\subsection{Selected Track Statistics}

\begin{center}
\begin{tabular}{@{}rlrrr@{}}
\toprule
Track ID & Branch & Samples & Pearson \(r\) & Spearman \(\rho\) \\
\midrule
3463 & left & 83 & 0.603518 & 0.561827 \\
254 & left & 83 & 0.756515 & 0.846252 \\
3127 & left & 84 & 0.803099 & 0.831943 \\
1713 & left & 84 & 0.837714 & 0.859534 \\
3460 & left & 75 & 0.527768 & 0.493072 \\
3461 & right & 67 & 0.668621 & 0.660548 \\
257 & right & 66 & 0.767774 & 0.779150 \\
3465 & right & 66 & 0.811301 & 0.717638 \\
3462 & right & 67 & 0.876987 & 0.876327 \\
1712 & right & 67 & 0.458025 & 0.526406 \\
\bottomrule
\end{tabular}
\end{center}

The pooled descriptive correlations were:

\[
r_{\mathrm{Pearson}} = 0.712629,
\qquad
\rho_{\mathrm{Spearman}} = 0.718781.
\]

\subsection{Outputs}

\begin{verbatim}
outputs/physics/video_2/droplet_branch_velocity_validation/selected_tracks.csv
outputs/physics/video_2/droplet_branch_velocity_validation/pooled_normalized_samples.csv
outputs/physics/video_2/droplet_branch_velocity_validation/droplet_branch_velocity_summary.json
outputs/physics/video_2/droplet_branch_velocity_validation/droplet_branch_velocity_scatter.png
outputs/physics/video_2/droplet_branch_velocity_validation/track_<track_id>_normalized_velocity.png
\end{verbatim}

\subsection{Execution}

\begin{verbatim}
.\.venv\Scripts\python.exe -m scripts.validate_droplet_branch_velocity \
  --experiment configs\experiments\video_2.yml \
  --overwrite
\end{verbatim}

\section{Generated Data Products}

\subsection{Visual Summary}

The figures in this section document the main visual artifacts generated during
the geometry, occupancy, baseline-hydraulics, and branch-velocity validation
steps. They are included directly from their generated repository locations so
the document continues to reference the same files consumed during validation.

\begin{figure}[H]
  \centering
  \includegraphics[width=0.92\linewidth,height=0.75\textheight,keepaspectratio]{../outputs/geometry/asymmetric_loop_h100/centerline_geometry.png}
  \caption{Centerline geometry diagnostic showing the ordered device centerline representation and derived branch structure.}
  \label{fig:centerline-geometry}
\end{figure}

\begin{figure}[H]
  \centering
  \includegraphics[width=0.95\linewidth,height=0.75\textheight,keepaspectratio]{../outputs/geometry/asymmetric_loop_h100/centerline_regions_frame_0.png}
  \caption{Legacy centerline-region diagnostic for frame 0. This plot was useful during development, but the current occupancy pipeline uses the device-level region labels shown in Figure~\ref{fig:region-overlay}.}
  \label{fig:centerline-regions-frame0}
\end{figure}

\begin{figure}[H]
  \centering
  \includegraphics[width=0.95\linewidth,height=0.75\textheight,keepaspectratio]{../outputs/physics/video_2/occupancy_diagnostics.png}
  \caption{Droplet fractional-occupancy diagnostics summarizing coverage, normalized occupancy behavior, and region assignment statistics.}
  \label{fig:occupancy-diagnostics}
\end{figure}

\begin{figure}[H]
  \centering
  \includegraphics[width=0.95\linewidth,height=0.75\textheight,keepaspectratio]{../outputs/physics/video_2/occupancy_spot_checks.png}
  \caption{Representative occupancy spot checks used to visually verify ellipse rasterization against the precomputed region labels.}
  \label{fig:occupancy-spot-checks}
\end{figure}

\begin{figure}[H]
  \centering
  \includegraphics[width=0.95\linewidth,height=0.75\textheight,keepaspectratio]{../outputs/physics/video_2/invalid_occupancy_random_10.png}
  \caption{Random sample of ten previously invalid or low-coverage occupancy cases, used for manual inspection of boundary and coverage behavior.}
  \label{fig:invalid-occupancy-random}
\end{figure}

\begin{figure}[H]
  \centering
  \includegraphics[width=0.95\linewidth,height=0.75\textheight,keepaspectratio]{../outputs/physics/video_2/baseline_hydraulic_diagnostics.png}
  \caption{Baseline branch-hydraulics diagnostics showing branch flow-rate and superficial-velocity behavior after normalized branch occupancy is converted into hydraulic resistance.}
  \label{fig:baseline-hydraulic-diagnostics}
\end{figure}

\begin{figure}[H]
  \centering
  \includegraphics[width=0.9\linewidth,height=0.75\textheight,keepaspectratio]{../outputs/physics/video_2/droplet_branch_velocity_validation/droplet_branch_velocity_scatter.png}
  \caption{Pooled normalized branch-velocity validation scatter. The x-axis is normalized calculated branch superficial velocity in \(\mu\mathrm{m}/\mathrm{s}\); the y-axis is normalized measured droplet speed derived in \texttt{px/frame}. The plot compares temporal variation after per-track normalization, not absolute velocity magnitudes.}
  \label{fig:droplet-branch-velocity-scatter}
\end{figure}

\begin{figure}[H]
  \centering
  \includegraphics[width=0.95\linewidth,height=0.72\textheight,keepaspectratio]{../outputs/physics/video_2/droplet_branch_velocity_validation/track_3463_normalized_velocity.png}
  \caption{Selected validation track 3463 in the left branch. The measured droplet-speed signal is derived by centered finite difference from \texttt{centroid\_x}/\texttt{centroid\_y}; both signals are shown as per-track z-scores.}
  \label{fig:validation-track-3463}
\end{figure}

\begin{figure}[H]
  \centering
  \includegraphics[width=0.95\linewidth,height=0.72\textheight,keepaspectratio]{../outputs/physics/video_2/droplet_branch_velocity_validation/track_254_normalized_velocity.png}
  \caption{Selected validation track 254 in the left branch. The comparison is zero-lag and unsmoothed.}
  \label{fig:validation-track-254}
\end{figure}

\begin{figure}[H]
  \centering
  \includegraphics[width=0.95\linewidth,height=0.72\textheight,keepaspectratio]{../outputs/physics/video_2/droplet_branch_velocity_validation/track_3127_normalized_velocity.png}
  \caption{Selected validation track 3127 in the left branch.}
  \label{fig:validation-track-3127}
\end{figure}

\begin{figure}[H]
  \centering
  \includegraphics[width=0.95\linewidth,height=0.72\textheight,keepaspectratio]{../outputs/physics/video_2/droplet_branch_velocity_validation/track_1713_normalized_velocity.png}
  \caption{Selected validation track 1713 in the left branch.}
  \label{fig:validation-track-1713}
\end{figure}

\begin{figure}[H]
  \centering
  \includegraphics[width=0.95\linewidth,height=0.72\textheight,keepaspectratio]{../outputs/physics/video_2/droplet_branch_velocity_validation/track_3460_normalized_velocity.png}
  \caption{Selected validation track 3460 in the left branch.}
  \label{fig:validation-track-3460}
\end{figure}

\begin{figure}[H]
  \centering
  \includegraphics[width=0.95\linewidth,height=0.72\textheight,keepaspectratio]{../outputs/physics/video_2/droplet_branch_velocity_validation/track_3461_normalized_velocity.png}
  \caption{Selected validation track 3461 in the right branch.}
  \label{fig:validation-track-3461}
\end{figure}

\begin{figure}[H]
  \centering
  \includegraphics[width=0.95\linewidth,height=0.72\textheight,keepaspectratio]{../outputs/physics/video_2/droplet_branch_velocity_validation/track_257_normalized_velocity.png}
  \caption{Selected validation track 257 in the right branch.}
  \label{fig:validation-track-257}
\end{figure}

\begin{figure}[H]
  \centering
  \includegraphics[width=0.95\linewidth,height=0.72\textheight,keepaspectratio]{../outputs/physics/video_2/droplet_branch_velocity_validation/track_3465_normalized_velocity.png}
  \caption{Selected validation track 3465 in the right branch.}
  \label{fig:validation-track-3465}
\end{figure}

\begin{figure}[H]
  \centering
  \includegraphics[width=0.95\linewidth,height=0.72\textheight,keepaspectratio]{../outputs/physics/video_2/droplet_branch_velocity_validation/track_3462_normalized_velocity.png}
  \caption{Selected validation track 3462 in the right branch.}
  \label{fig:validation-track-3462}
\end{figure}

\begin{figure}[H]
  \centering
  \includegraphics[width=0.95\linewidth,height=0.72\textheight,keepaspectratio]{../outputs/physics/video_2/droplet_branch_velocity_validation/track_1712_normalized_velocity.png}
  \caption{Selected validation track 1712 in the right branch.}
  \label{fig:validation-track-1712}
\end{figure}

\subsection{Geometry Products}

\begin{itemize}
  \item \texttt{outputs/geometry/asymmetric\_loop\_h100/centerline\_geometry.csv}
  \item \texttt{outputs/geometry/asymmetric\_loop\_h100/geometry\_summary.json}
  \item \texttt{outputs/geometry/asymmetric\_loop\_h100/device\_geometry.npz}
  \item \texttt{outputs/geometry/asymmetric\_loop\_h100/centerline\_geometry.png}
\end{itemize}

\subsection{Device Region Products}

\begin{itemize}
  \item \texttt{data/geometry/asymmetric\_loop\_h100/region\_labels.npy}
  \item \texttt{data/geometry/asymmetric\_loop\_h100/region\_metadata.json}
  \item \texttt{data/geometry/asymmetric\_loop\_h100/region\_overlay.png}
\end{itemize}

\subsection{Occupancy Products}

\begin{itemize}
  \item \texttt{outputs/physics/video\_2/droplet\_occupancy.csv}
  \item \texttt{outputs/physics/video\_2/droplet\_occupancy\_summary.json}
  \item \texttt{outputs/physics/video\_2/occupancy\_diagnostics.png}
  \item \texttt{outputs/physics/video\_2/occupancy\_spot\_checks.png}
\end{itemize}

\subsection{Hydraulic Products}

\begin{itemize}
  \item \texttt{outputs/physics/video\_2/baseline\_hydraulic\_state.csv}
  \item \texttt{outputs/physics/video\_2/baseline\_hydraulic\_summary.json}
  \item \texttt{outputs/physics/video\_2/baseline\_hydraulic\_diagnostics.png}
\end{itemize}

\subsection{Branch-Velocity Validation Products}

\begin{itemize}
  \item \texttt{outputs/physics/video\_2/droplet\_branch\_velocity\_validation/selected\_tracks.csv}
  \item \texttt{outputs/physics/video\_2/droplet\_branch\_velocity\_validation/pooled\_normalized\_samples.csv}
  \item \texttt{outputs/physics/video\_2/droplet\_branch\_velocity\_validation/droplet\_branch\_velocity\_summary.json}
  \item \texttt{outputs/physics/video\_2/droplet\_branch\_velocity\_validation/droplet\_branch\_velocity\_scatter.png}
  \item \texttt{outputs/physics/video\_2/droplet\_branch\_velocity\_validation/track\_<track\_id>\_normalized\_velocity.png}
\end{itemize}

\section{Testing}

The test suite currently covers:

\begin{itemize}
  \item centerline arc length, orientation, tangent, normal, and topology behavior;
  \item configuration-to-device resolution;
  \item region label validity and connectivity;
  \item ellipse rasterization and occupancy normalization;
  \item low-coverage diagnostics and computability semantics;
  \item baseline hydraulic equations, unit conversion, flow conservation, and mixture-flow handling;
  \item centered finite-difference velocity derivation, branch-interior filtering,
  deterministic candidate selection, correlation calculations, and validation metadata.
\end{itemize}

Run all tests with:

\begin{verbatim}
.\.venv\Scripts\python.exe -m pytest tests
\end{verbatim}

\section{Scope Boundaries}

The current implementation does not:

\begin{itemize}
  \item modify files under \texttt{data/processed/2};
  \item train or modify a neural network;
  \item construct transformer training windows;
  \item model junction resistance;
  \item model droplet--droplet interactions;
  \item fit hydraulic parameters from data;
  \item use segmented droplet masks;
  \item calculate phase-specific slip or separate phase velocities;
  \item prove that branch superficial velocity equals droplet speed;
  \item optimize temporal lag or smooth the validation signals.
\end{itemize}

\section{Recommended Pipeline Order}

\begin{enumerate}[label=\arabic*.]
  \item Build centerline geometry:
\begin{verbatim}
.\.venv\Scripts\python.exe -m scripts.build_geometry \
  --experiment configs\experiments\video_2.yml \
  --overwrite
\end{verbatim}

  \item Build device region labels:
\begin{verbatim}
.\.venv\Scripts\python.exe -m scripts.build_device_regions \
  --experiment configs\experiments\video_2.yml \
  --frame-index 0 \
  --overwrite
\end{verbatim}

  \item Compute droplet occupancy:
\begin{verbatim}
.\.venv\Scripts\python.exe -m scripts.compute_droplet_occupancy \
  --experiment configs\experiments\video_2.yml \
  --overwrite
\end{verbatim}

  \item Compute baseline hydraulics:
\begin{verbatim}
.\.venv\Scripts\python.exe -m scripts.compute_baseline_hydraulics \
  --experiment configs\experiments\video_2.yml \
  --overwrite
\end{verbatim}

  \item Run droplet branch-velocity validation:
\begin{verbatim}
.\.venv\Scripts\python.exe -m scripts.validate_droplet_branch_velocity \
  --experiment configs\experiments\video_2.yml \
  --overwrite
\end{verbatim}
\end{enumerate}

\end{document}
